\documentclass[11pt,letterpaper]{article}
\usepackage[T1]{fontenc}
\usepackage{lmodern,amsmath,amssymb,amsthm,microtype}
\usepackage[margin=1in]{geometry}
\usepackage[hidelinks]{hyperref}
\providecommand{\tightlist}{\setlength{\itemsep}{0pt}\setlength{\parskip}{0pt}}
\DeclareUnicodeCharacter{2264}{\ensuremath{\le}}
\DeclareUnicodeCharacter{2265}{\ensuremath{\ge}}
\DeclareUnicodeCharacter{2260}{\ensuremath{\ne}}
\DeclareUnicodeCharacter{2212}{\ensuremath{-}}

\DeclareUnicodeCharacter{2211}{\ensuremath{\sum}}

\usepackage{longtable,booktabs,array}
\title{Conditional XYZ and a two-arm case of the balanced-pair conjecture}
\author{Aafi Khan Malek\\ \small Independent researcher, Ahmedabad, India\\
\small\href{mailto:aafimalek2032@gmail.com}{aafimalek2032@gmail.com}}
\date{28 November 2026}
\begin{document}
\maketitle

Let \(P\) be a finite poset, \(a,b\) incomparable, and suppose each
strict predecessor of \(a\) is below \(b\). Let
\(S={b}\cup(U(b)\setminus U(a))\), where \(U\) means strict successors.
Each element of \(S\) is incomparable with \(a\). All probabilities are
for a uniformly random linear extension, and \(F(z)=\Pr(a<z)\) uses
extension comparison. Fix \(p=F(b)\le1/2\). A balanced pair is any
distinct pair with comparison probability in \([1/3,2/3]\). In the
analysis by contradiction, assume none; then \(p<1/3\), every vertex
\(z\in S\) is low \((F(z)<1/3)\) or high \((F(z)>2/3)\), and \(b\) is
low.

\section{Two known rank facts, for
reference}\label{two-known-rank-facts-for-reference}

Fix an induced order \(R\) of \(S\), numbering ranks from one. Let
\(m_R(k)\) be the unconditional probability of induced order \(R\) and
of \(a\) coming after exactly \(k\) points of \(S\). This sequence is
antitone in \(k\): for \(k>0\) swap \(a\) with the last element \(z\) of
\(S\) before \(a\). Moving \(a\) earlier crosses no predecessor of
\(a\), as every predecessor lies below \(b\le z\). Any successor of
\(z\) strictly between \(z\) and \(a\) would lie above \(b\) and not
above \(a\), hence belong to \(S\), contradiction. The swap preserves
\(R\) and can be inverted by swapping \(a\) with the next \(S\)-element.
Summing bins over orders gives an antitone total sequence \(q_k\), with
\(q_0=p\) and \(q_{|S|}\le p\).

Write \(M_R(h)=\sum_{k<h}m_R(k)\). Antitonicity gives both
\(M_R(h+k)\le M_R(h)+M_R(k)\) (truncate at the final bin if needed) and
\(M_R(v)/v\le M_R(u)/u\) for positive \(u\le v\). Thus rank domination
\(\rho_R(z)\le\rho_R(u)+\rho_R(v)\) in every \(R\) implies
\(F(z)\le F(u)+F(v)\). Repeated anchors are allowed. Also if
\(\rho_R(z)\le c\rho_R(u)\) in every \(R\) and \(c\ge1\), then
\(F(z)\le cF(u)\), by the average inequality and monotonicity.

\section{Conditional XYZ inequality for two high
successors}\label{conditional-xyz-inequality-for-two-high-successors}

Let a low \(q\in S\) precede distinct \(x,y\in S\) in the poset. Write
\(s=F(q)\), \(u=\Pr(x<a)\), \(v=\Pr(y<a)\), \(J=\Pr(x<a,y<a)\), and
\(B=\Pr(a<x,a<y)\). Both events defining the marginals \(u,v\) imply
\(q<a\). The conditional law given \(q<a\) is uniform on extensions of
the poset with \(q<a\) adjoined and closed transitively. Apply Shepp's
XYZ inequality to the two same-sided comparisons with \(a\), using the
dual poset if needed. Clearing denominators gives \[
   uv\le (1-s) J=(1-s)(u+v-1+B).
\] If \(x,y\) are both high then \(0\le u,v<1/3\). On the square
\([0,1/3]^2\), the function \(G(u,v)=(1-s)(u+v-1+B)-uv\) is strictly
increasing in each variable, since \(1-s>2/3\). Consequently \[
  \frac19 < (1-s)(B-\tfrac13).                    \tag{1}
\] In particular, \textbf{both cannot be high if} \[
  2\min(\rho_R(x),\rho_R(y))\le3\rho_R(q)
  \quad\text{for every }R.                         \tag{2}
\] Indeed, \(B=\sum_R M_R(\min(\rho_R(x),\rho_R(y)))\le\frac32s\) by the
antitone prefix average inequality. For \(0\le s<1/3\), \[
  \tfrac19-(1-s)(\tfrac32s-\tfrac13)
  =(\tfrac13-s)(\tfrac43-\tfrac32s)>0,
\] contradicting (1).

For example, if \(S\) starts with an ordinal prefix of \(h\ge2\)
vertices ending in a low maximum \(q\), the rank of \(q\) is always
\(h\). If the remainder has exactly two minimal elements \(x,y\), their
smaller rank is \(h+1\), so (2) shows at least one is low. This is a
local step: it does not assume both high vertices themselves have
bounded ranks.

\section{Infinite two-arm class at the original good-pair
cap}\label{infinite-two-arm-class-at-the-original-good-pair-cap}

\textbf{Theorem.} A balanced pair exists if \(S\) is an ordinal sum of a
chain of \(h\ge2\) vertices starting at \(b\) and a tail partitioned
into two nonempty chains \(x_1<\dots<x_r\), \(y_1<\dots<y_t\) with
\(r,t\le h+1\). Additional relations between the tail chains are
allowed. The original branch hypotheses above, including \(p\le1/2\),
suffice.

In a putative no-balance case, all prefix-chain elements are low: at
rank \(i\ge2\), the point has rank at most twice the previous rank
\(i-1\). If the two heads \(x_1,y_1\) are incomparable, they are both
minimal in the tail and the conditional step shows at least one is low.
If comparable, the smaller has rank \(h+1\le2h\) and is low by pooling
the last prefix vertex with itself. The other head has rank at most
\(h+(h+1)+1=2h+2\) and each head has rank at least \(h+1\), so pooling
the low head with itself makes both low.

A first high vertex at \(x_2\) would have rank at most \(2h+3\). The two
distinct low heads always have rank sum at least \(h+1+h+2=2h+3\),
contradiction. For a first high vertex \(x_{i+1}\) with \(i\ge2\), its
rank is at most \(h+i+1+t\le2h+i+2\le2(h+i)\le2\rho_R(x_i)\), again
impossible. The argument for the other chain is identical. Therefore
both tail ends are low. Every extension either places \(a\) before at
least one tail end or puts every point of \(S\) before \(a\). Its
probability one is consequently at most \(F(x_r)+F(y_t)+q_{|S|}<1\),
contradiction.

\section{Consequences at the bottom of the branch and sharp histogram
boundary}\label{consequences-at-the-bottom-of-the-branch-and-sharp-histogram-boundary}

If the two nonroot minima \(x,y\) of an exclusive branch are both high,
then \(\min(\rho_R(x),\rho_R(y))=2\). Taking \(q=b\) gives
\(B=q_0+q_1\le2p\). Inequality (1) now forces \[
   (1-p)(2p-\tfrac13)>\tfrac19,
   \qquad p>\frac{7-\sqrt {17}}{12}=0.23974\ldots.  \tag{3}
\] This bound does not depend on the tail lengths, but it requires that
both nonroot minima are high. It also applies to any branch we can prove
has that property under no balance.

In particular, for a chain \(x_1<\dots<x_r\) plus one vertex \(y\) above
\(b\), under no balance all nonroot vertices are high. To see this
directly, both tail ends \(x_r,y\) must be high: if both were low, the
two-end cover and final-bin bound give a contradiction. If only one were
low, consider the first high of the other chain. If the low chain has
length \(t\) and the other chain has \(i\) low vertices before its first
high, that first high has rank at most \(i+t+2\). The later of the two
low anchors (each chain's last low, or \(b\) if it has none) has rank at
least \(i+t+1\), because both low prefixes have appeared by then. Adding
the earlier anchor, of rank at least one, gives the desired rank sum.
This contradicts the pooled-rank bound. Now, if \(x_1\) were low, let
\(x_{i+1}\), \(i\ge1\), be the first high on the long chain. Its rank is
at most \(i+3\le2(i+1)\le2\rho_R(x_i)\), contradiction.

For the incomparable-minima chain-plus-one case let \(d\ge1\) be the
number of long-chain vertices not strictly above \(y\). In every branch
order, \[
   d\rho_R(x_1)+\rho_R(y)\le 3d+2.
\] If \(y\) is first after \(b\), the two ranks are \(3,2\); otherwise
\(x_1\) is second and \(\rho_R(y)\le d+2\). Since
\(M_R(k)\le k\,m_R(0)\), averaging gives \[
 d F(x_1)+F(y)\le (3d+2)p.
\] Thus for \(d\ge2\) the chain-plus-one class has balance at \[
 p\le\max\left\{
    \frac{2(d+1)}{3(3d+2)},\frac{7-\sqrt{17}}{12}
 \right\}.
\] The maxima are \(1/4\) for \(d=2\), \(8/33\) for \(d=3\), and
\((7-\sqrt{17})/12\) for every \(d\ge4\). When \(d=1\), the earlier
\(5/18\) bound is stronger: let \(z\) be the majority winner between the
two heads \(x_1,y\), and let \(w\) be the other. In any fixed order
normalized to mass one, write \(P_0\) for its cut-zero probability and
\(A_z\) for its probability of \(a<z\). If \(z<w\), then
\(A_z/2\le P_0\) because the rank of \(z\) is two; otherwise the rank is
at most three, so \(A_z/2-1/6\le P_0\) by \(A_z\le\min(1,3P_0)\).
Averaging gives \(p\ge F(z)/2-(1-\Pr(z<w))/6>5/18\) under no balance.
Comparable minima already give balance at \(p\le1/2\), since the smaller
one has rank two. The weighted rank argument also applies when both
heads of arbitrary two arms are high: putting
\(d_x=\#\{x_i:y_1\not<x_i\}\) and \(d_y=\#\{y_j:x_1\not<y_j\}\), it
gives \[
  d_xF(x_1)+d_yF(y_1)\le(2d_x+2d_y+d_xd_y)p.
\]

The weighted chain-plus-one bound is sharp for the \textbf{antitone
histogram abstraction} when \(d\ge2\). Take \(r=d\) with no additional
relations. A component of mass \(w\) has \(h'\) equally likely cuts
\(0,\ldots,h'-1\) of one branch order. Use these three components:

\begin{longtable}[]{@{}lrr@{}}
\toprule\noalign{}
Branch order & Number \(h'\) of cuts & Mass \\
\midrule\noalign{}
\endhead
\bottomrule\noalign{}
\endlastfoot
\(b,y,x_1,\ldots,x_d\) & 3 & \((d-2)/(3d+2)\) \\
\(b,y,x_1,\ldots,x_d\) & 4 & \(4(d+2)/(3(3d+2))\) \\
\(b,x_1,\ldots,x_d,y\) & \(d+2\) & \(2(d+2)/(3(3d+2))\) \\
\end{longtable}

Then \(p=2(d+1)/(3(3d+2))\), \(F(x_1)=F(y)=2/3\), and
\(\Pr(y<x_1)\ge2/3\). Mix an arbitrarily small positive mass with the
first order at the only cut 0. This makes all three relations strict,
and preserves \(p<1/3\), exhibiting unbalanced histograms arbitrarily
close above the threshold. These are \textbf{not} counterexamples to the
poset theorem. At the threshold \(J=2p-1/3=(d+2)/(3(3d+2))\), and for
\(d\ge4\), \(uv=1/9>(1-p)J\); after clearing the denominator the sign is
that of \(2d^2-6d-4>0\). Conditional XYZ, absent from the histogram
abstraction, rejects these exact instances and their sufficiently small
perturbations.

Reference: L. A. Shepp, \emph{The XYZ Conjecture and the FKG
Inequality}, Ann. Probab. 10 (1982), 824--827,
\href{https://doi.org/10.1214/aop/1176993791}{doi:10.1214/aop/1176993791}.

Scope: these proofs are written mathematics. The conditional XYZ use and
new branch theorem have not been checked in Lean. The full width-three
conjecture remains open.

\end{document}
